
\begin{remark}\label{fga-rem-3-2}\dcref{3.2}\source{fga-explained:page-0051}\uses{}
\begin{verbatim}
The definitionof cartesianarrow we give is more restrictive
than the definition
                  in [SGAl]; our cartesianarrows are calledstronglycartesianin
[Gra66]. However, the resultingnotionsof fiberedcategorycoincide.
42                           3.FIBERED CATEGORIES
\end{verbatim}
\end{remark}


\begin{remark}\label{fga-rem-3-3}\dcref{3.3}\source{fga-explained:page-0052}\uses{}
\begin{verbatim}
Given two pullbacks(p-?,^ r]and (/>:
                                                   ^ ^ ?7of 77to [/,the
unique arrow 0: ^ ^^ ^ that fits
                               intothe diagram




                                              >F
isan isomorphism; the inverseisthe arrow ^ ^ ^ obtained by exchanging ^ and ^
in the diagram above.
    In other words, a pullbackisunique,up to a unique isomorphism.
     The followingfactsare easy to prove,and are leftto the reader.
\end{verbatim}
\end{remark}


\begin{proposition}\label{fga-pro-3-4}\dcref{3.4}\source{fga-explained:page-0052}\uses{}
\begin{verbatim}
(i)If^ isa categoryover C, the composite of cartesianarrows in T iscartesian,
  (ii)An arrow in T whose image in C isan isomorphism iscartesianifand only if
      itisan isomorphism,
 (iii)If ^ ^ ?7and rj^ ( are arrows in T and rj^ ( is cartesian,then ^ ^ 77is
      cartesianifand only ifthe composite ^ ^ ^ iscartesian,
 (iv)Let pg: G ^ C and F: T ^^ G he functors,(/>:   ^ ^ 77an arrow in T. If (pis
      cartesianoveritsimage Fcj):F^ -^ Frj in Q and Fcj)iscartesianoveritsimage
      PgFcp: pgF^ -^ PtFt] in C, then (piscartesianover itsimage pgFcp in C.
\end{verbatim}
\end{proposition}
\begin{proof}
\notready
\end{proof}


\begin{definition}\label{fga-def-3-6}\dcref{3.6}\source{fga-explained:page-0052}\uses{}
\begin{verbatim}
If^ and Q are fiberedcategoriesover C, then a morphism of
fiberedcategoriesF: T ^ G is8ifunctorsuch that:
                         that is,Pg 0 F = pj^;
  (i)F is base-preserving,
 (ii)F  sends cartesian
                      arrows to cartesianarrows.
    Noticethatinthe definition above the equalitypgoF = pjrmust be interpreted
as an actualequality.In other words, the existenceof an isomorphism of functors
between pg o F and pjr isnot enough.
\end{verbatim}
\end{definition}


\begin{definition}\label{fga-def-3-10}\dcref{3.10}\source{fga-explained:page-0054}\uses{}
\begin{verbatim}
A pseudo-functor$ on C consistsof the followingdata.
   (i)For each object [/of C a category ^U.
  (ii)For each arrow /: U ^ V a,functor/* :^V -^ ^U.
 (iii)For each object[/ofC an isomorphism eu' id^ :^ id^u offunctors^U -^ ^U.
 (iv)For each pair of arrows U -^V ^W     mi isomorphism

                      a/,,:/V-(^/)*:^W        ?^f^
    of functors^W -^ ^U.
   These data are requiredto satisfythe followingconditions.
(a) /: [/ ^ F isan arrow in C and 77isan objectof ^V, we have
   If
                              = ^u{f*v)- idL//> ?' rv
                      (^idujijl)
     and


(b) Whenever we have arrows [/ -^ F ^ Ty -^ r and an object 9 of J^(r),the
    diagram
                       f^g^h^e ????>         {gfYh'^e
                           rocgM^)             o^ghjiO)
                                 "^'""^'^
                       rihgre             Hhg/re
     commutes.
    In thisdefinition  we only consider(contravariant) pseudo-functorintothe
     of
category           O f
        categories. course,    thereisa much  more generalnotionof pseudo-functor
with valuesin a 2-category,which we willnot use at all.
    A functor $ from S into the category of categoriescan be consideredas a
pseudo-functor,in which every eu is the identityon $[/, and every af^g is the
identityon/*^* = (^/)*.
    We have seen how to associatewith a fiberedcategoryover C, equipped with a
cleavage,the data fora pseudo-functor;  we still
                                               have to checkthatthetwo conditions
of the definition
                are satisfied.
\end{verbatim}
\end{definition}


\begin{remark}\label{fga-rem-3-27}\dcref{3.27}\source{fga-explained:page-0064}\uses{}
\begin{verbatim}
It is interestingto noticethat ifF\ C?p -^ (Set)isa functor
and T ^^ C the associatedcategory fiberedin sets,then an object (X,^) of ^ is
universalpair forthe functorF ifand only ifitisa terminalobject forT. Hence
F isrepresentableifand only ifT has a terminalobject.

                  given an objectX of C, we have the representablefunctor
     In particular,

                               hx:C?P ?(Set),
definedon objectsby the rule h^t^ = Homc(t/,X). The category in setsover C
associatedwith thisfunctoristhe comma category(C/X), and the functor(CjX) -^
C isthe functorthat forgetsthe arrow intoX.
    So the situationisthe following.Prom Yoneda's lemma we seethatthe category
C isembedded intothe category of functorsC?p -^ (Set),while the category of
functorsisembedded intothe categoryof fiberedcategories.
    Prom now we willidentifya functorF\ C?p -^ (Set)with the corresponding
categoryfiberedin setsover C, and will(inconsistently) calla category fiberedin
setssimply "a functor".

     3.4.1. Categories fibered over an object.
\end{verbatim}
\end{remark}


\begin{proposition}\label{fga-pro-3-28}\dcref{3.28}\source{fga-explained:page-0064}\uses{}
\begin{verbatim}
Let 5 be a category fiberedin setsover C, T another
category,F\ T ^ Q o,functor.Then T isfiberedover Q ifand only ifitisfibered
over C via the composite pg o F: ^ ^ C.
    Furthermore, T is fiberedin groupoids over Q ifand only ifit fiberedin
groupoids over C, and is fiberedin setsover Q ifand only ifit fiberedin sets
over C.
3.4.FUNCTORS AND CATEGORIES FIBERED IN SETS                  55
\end{verbatim}
\end{proposition}
\begin{proof}
\begin{verbatim}
One sees immediately that an arrow of Q iscartesianover itsimage
in T ifand only ifitiscartesianover itsimage in C, and the firststatement follows
from this.
    Furthermore, one seesthat the fiberof T over an object [/of C isthe disjoint
union,as a category,of the fibersof T over allthe objectsof Q over U; thesefiber
are groupoids,or sets,ifand only iftheirdisjointunion is.                     D

    This can be used as follows.Suppose that S isan object of C, and consider
the categoryfiberedin sets(C/S) -^ C, correspondingto the representablefunctor
h^: C?P -^ (Set).By Proposition3.28,a fiberedcategoryT -^ (C/S) isthe same
as a fiberedcategory T ^ C, togetherwith a morphism T -^ {C/S) of categories
fiberedover C.
    In particular,categoriesfiberedinsetscorrespondto functors;hence we get that
giving a functor (C/Sy^ -^ (Set)is equivalentto assigninga functorC?p -^ (Set)
together with a naturaltransformation F ^ h^. Describingthisprocessforfunctors
seems lessnaturalthan forfiberedcategoriesin general.
    Given a functorF\ (C/S)^^ -^ (Set),this corresponds to a category fibered
in setsF -^ {C/S), which can be composed with the forgetfulfunctor {C/S) -^ C
to get a category fiberedin setsF ^ C, which in turn corresponds to a functor
F': C?P -^ (Set).What isthisfunctor? One minute's thought willconvince you
that itcan be describedas follows:F'{U) isthe disjointunion of the F{U -^ S)
forallthe arrows [/ ^ 5 in C. The actionof F' on arrows isthe obvious one.

   3.4.2. Fibered subcategories.
\end{verbatim}
\end{proof}


\begin{definition}\label{fga-def-3-34}\dcref{3.34}\source{fga-explained:page-0066}\uses{}
\begin{verbatim}
Let T and Q be two fiberedcategoriesover C. An
       of T with 5 is a morphism F\ T ^ Q, such that there existsanother
equivalence,
morphism G: Q -^ T, togetherwith isomorphisms of G o F with idjrand of F o G
with idg.
    We callG simply an inverseto F.
\end{verbatim}
\end{definition}


\begin{proposition}[fiberwise equivalence]\label{fga-pro-3-36}\dcref{3.36}\source{fga-explained:page-0066}\uses{}
Let $F:\mathcal T\to\mathcal G$ be a morphism of fibered categories over
$\mathcal C$.  Then $F$ is an equivalence if and only if, for every object
$U$ of $\mathcal C$, the induced functor
$F_U:\mathcal T(U)\to\mathcal G(U)$ is an equivalence.
\end{proposition}
\begin{proof}
An inverse of $F$ restricts to an inverse of every $F_U$.  Conversely, assume
that each $F_U$ is an equivalence.  The fiberwise full-faithfulness lemma below
shows that $F$ is fully faithful.  Choose, for every object $\eta$ of
$\mathcal G$, an object $G\eta$ in the fiber over $p(\eta)$ and an isomorphism
$a_\eta:F(G\eta)\xrightarrow{\sim}\eta$.  Full faithfulness gives a unique
arrow $Gf:G\eta\to G\eta'$ over every arrow $f:\eta\to\eta'$ whose image is
$a_{\eta'}^{-1}fa_\eta$.  These arrows define a functor $G$, and $a$ is an
isomorphism $FG\simeq\operatorname{id}$.  Applying full faithfulness to
$a_{F\xi}$ gives an isomorphism $GF\simeq\operatorname{id}$, so $F$ is an
equivalence.
\end{proof}


\begin{lemma}[fiberwise full faithfulness]\label{fga-lem-3-37}\dcref{3.37}\source{fga-explained:page-0066}\uses{}
Let $F:\mathcal T\to\mathcal G$ be a morphism of fibered categories such that
every $F_U:\mathcal T(U)\to\mathcal G(U)$ is fully faithful.  Then $F$ is
fully faithful.
\end{lemma}
\begin{proof}
Let $\xi,\eta$ be objects of $\mathcal T$ and let $f:F\xi\to F\eta$
lie over $u:p(\xi)\to p(\eta)$.  Choose a cartesian lift
$\bar\eta\to\eta$ of $u$ in $\mathcal T$.  Its image is a cartesian lift of
$u$ in $\mathcal G$.  The arrow $f$ therefore factors uniquely through
$F\xi\to F\bar\eta$ by an arrow in the fiber over $p(\xi)$.  Fiberwise full
faithfulness lifts that arrow uniquely to $\xi\to\bar\eta$, and composition
with $\bar\eta\to\eta$ gives the required unique lift of $f$.
\end{proof}


\begin{proposition}[categories equivalent to sets]\label{fga-pro-3-38}\dcref{3.38}\source{fga-explained:page-0068}\uses{}
A category is equivalent to a set, viewed as a discrete category, if and only
if it is an equivalence relation: between any two objects there is at most one
arrow, every arrow is invertible, and every object has its identity arrow.
\end{proposition}
\begin{proof}
If a category is equivalent to a set, it has precisely this property.  Conversely,
let $(X,R)$ be an equivalence relation and let $X/R$ be the set of equivalence
classes.  The functor $(X,R)\to X/R$ sending an object to its class is fully
faithful and essentially surjective, hence an equivalence.
\end{proof}


\begin{definition}[quasi-functor]\label{fga-def-3-39}\dcref{3.39}\source{fga-explained:page-0068}\uses{fga-def-3-1}
A category $\mathcal F$ over $\mathcal C$ is a quasi-functor, or is fibered in
equivalence relations, if it is fibered and each fiber $\mathcal F(U)$ is an
equivalence relation.
\end{definition}


\begin{proposition}[quasi-functor criterion]\label{fga-pro-3-40}\dcref{3.40}\source{fga-explained:page-0068}\uses{fga-def-3-39}
A category $\mathcal F$ over $\mathcal C$ is a quasi-functor if and only if the
following two conditions hold.
\begin{enumerate}
\item Given an object $\eta$ of $\mathcal F$ and an arrow
$f:U\to p_{\mathcal F}\eta$ of $\mathcal C$, there exists an arrow
$\phi:\xi\to\eta$ of $\mathcal F$ with $p_{\mathcal F}\phi=f$. Furthermore,
given any other arrow $\phi':\xi'\to\eta$ with
$p_{\mathcal F}\phi'=f$, there exists $\alpha:\xi\to\xi'$ in $\mathcal F(U)$
such that $\phi'\alpha=\phi$.
\item Given two objects $\xi$ and $\eta$ of $\mathcal F$ and an arrow
$f:p_{\mathcal F}\xi\to p_{\mathcal F}\eta$ of $\mathcal C$, there exists at
most one arrow $\xi\to\eta$ over $f$.
\end{enumerate}
\end{proposition}
\begin{proof}
\notready
\end{proof}


\begin{proposition}[quasi-functors and functors]\label{fga-pro-3-41}\dcref{3.41}\source{fga-explained:page-0068}\uses{fga-pro-3-36,fga-pro-3-38,fga-def-3-39,fga-prop-3-22}
A fibered category over $\mathcal C$ is a quasi-functor if and only if it is
equivalent to a functor.
\end{proposition}
\begin{proof}
This is an application of Proposition 3.36.  Suppose first that a
fibered category $\mathcal F$ is equivalent to a functor $\Phi$.  Then every
category $\mathcal F(U)$ is equivalent to the set $\Phi(U)$, so $\mathcal F$ is
fibered in equivalence relations over $\mathcal C$ by Proposition 3.38.

Conversely, assume that $\mathcal F$ is fibered in equivalence relations.  In
particular it is fibered in groupoids, so every arrow in $\mathcal F$ is
cartesian by Proposition 3.22.  For each object $U$ of
$\mathcal C$, let $\Phi(U)$ be the set of isomorphism classes of objects of
$\mathcal F(U)$.  Given $f:U\to V$ in $\mathcal C$, pullback defines
$f^*:\Phi(V)\to\Phi(U)$: isomorphic objects and their pullbacks remain
isomorphic, so this is well defined.  These maps give a functor
$\Phi:\mathcal C^{\mathrm{op}}\to\mathbf{Set}$.

Regard $\Phi$ as a category fibered in sets.  The canonical morphism
$\mathcal F\to\Phi$ is an equivalence on each fiber, and hence is an
equivalence by Proposition 3.36.
\end{proof}


\begin{proposition}[morphisms into a fibered category]\label{fga-pro-3-42}\dcref{3.42}\source{fga-explained:page-0068}\uses{fga-def-3-39}
Let $\mathcal F$ and $\mathcal G$ be fibered categories over $\mathcal C$.
\begin{enumerate}
\item If $\mathcal G$ is fibered in groupoids, then
$\operatorname{Hom}_{\mathcal C}(\mathcal F,\mathcal G)$ is a groupoid.
\item If $\mathcal G$ is a quasi-functor, then
$\operatorname{Hom}_{\mathcal C}(\mathcal F,\mathcal G)$ is an equivalence
relation.
\item If $\mathcal G$ is a functor, then
$\operatorname{Hom}_{\mathcal C}(\mathcal F,\mathcal G)$ is a set.
\end{enumerate}
\end{proposition}
\begin{proof}
\notready
\end{proof}


\begin{definition}\label{fga-def-3-43}\dcref{3.43}\source{fga-explained:page-0069}\uses{}
\begin{verbatim}
A fiberedcategoryover C is representable
                                                           ifitisequivalent
to a categoryof the form {C/X).

    So a representablecategoryisnecessarily
                                          a quasi-functor,
                                                         by Proposition3.41.
However, we should be careful:ifT and Q are fiberedcategories, equivalentto
{C/X) and {C/Y) fortwo objectsX and Y of C, then
                      Hom(X,y)       = Hom((C/X), {C/Y)),

and accordingto Proposition3.35we have an equivalenceof categories

                     Hom((C/X), {C/Y)) 0^ Homc(^, G);
but Homc(^, G) need not be a set,itcould very wellbe an equivalencerelation.

    3.6.2. The 2-categorical Yoneda lemma. As in the case of functors,we
have a strongerversionof the 2-categorical Yoneda lemma. Suppose that ^ isa
categoryfiberedover C, and thatX  isan objectofC. Let therebe given a morphism
F: {C/X) -^ T\ with thiswe can associatean objectF(idx) G ?^(-^)-Also,to each
base-preservingnaturaltransformationa: F ^ G of functorsF, G: {C/X) -^ T
we associatethe arrow aid^ ? ^(idx) -^ G(idx). This definesa functor

                         Home ((C/X),^) ?^T{X).

    Conversely,given an object ^ G T{X) we get a functor F^: {C/X) ^ ^ as
follows.Given an object4>\U ^Xoi    (C/X), we defineF^{4>)= (/)*^
                                                                G T{U); with
an arrow
                                 u            yv


                                         X
60                          3.FIBERED CATEGORIES

                                           -^ ip*^in ^{17) making the diagram
in {C/X) we associatethe only arrow 9\ cj)*^




commutative. We leaveitto the readerto check that F^ isindeed a functor.

     2-YoNEDA   Lemma.   The two functorsabove definean equivalenceof categories
                           Home ((C/X),^):^^(X).
\end{verbatim}
\end{definition}
\begin{proof}
\begin{verbatim}
To check that the composite

                     T{X) ?. Home ((C/X),^) ? ^(X)
isisomorphicto the identity,noticethat forany object^ G ^(X), the composition
        to               =
applied ^ yieldsF^{^) id^^:^,    which is canonicallyisomorphic to ^. It is easy
to see that thisdefinesan isomorphism of functors.
    For the composite
                Home ((C/X),^) ?. T{X) ?^ Home ((C/X),^)
take a morphism F: (C/X) -^ T and set ^ = F(idx). We need to produce a base-
preservingisomorphism of functorsof F with F^. The identityidx isa terminal
objectin the category{C/X), hence forany object(p:U ^ X thereisa unique
arrow
    (p:idx, which isclearlycartesian.Hence itwillremain cartesianafterapplying
F, because F isa functor:thismeans that F^cf)) isa pullbackof ^ = F(idx) along
(j): ^
   [/ X, so thereisa canonicalisomorphism F^((/))   = (/)*?,
                                                          ? F{(j))in T{U). Itis
easy to check that                                         of
                   thisdefines base-preservingisomorphism functors,and this
                              a
ends the proof.                                                              D

    We have identifiedan objectX with the functorhx '-C?p -^ (Set)itrepresents,
and we have identifiedthe functorhx with the correspondingcategory(C/X): so,
to be consistent,
                we have to identifyX and (C/X). So, we willwriteX for (C/X).
    As forfunctors,the strong form of the 2-Yoneda Lemma can be used to
reformulate
       the conditionof representability.A morphism (C/X) -^ T correspondsto
an object^ G T{X), which in turn definesthe functorF' \ {C/X) -^ T described
above; thisis isomorphic to the originalfunctorF. Then F' isan equivalenceif
and only ifforeach object IJoiC the restriction

                   F[j'.Home(/7,X) = {C/X){U) ?^ T{V)
that sends each f \ U -^ X to the pullback /*^ G TiJJ), is an equivalenceof
categories.Since Home(t/, X) isa set,thisisequivalentto saying that TiJJ) isa
groupoid,and each objectof TiJJ) isisomorphicto the image of a unique element
of Home(t/,X) via a unique isomorphism. Sincethe isomorphisms p c^ /*^ in U
correspond to cartesianarrows p ^ ^, and in a groupoid allarrows are cartesian,
thismeans that T isfiberedin groupoids,and foreach p G TiJJ) there existsa
unique arrow p ^ ^. We have proved the following.
3.7.THE FUNCTORS OF ARROWS OF A FIBERED CATEGORY                  61
\end{verbatim}
\end{proof}


\begin{proposition}\label{fga-pro-3-44}\dcref{3.44}\source{fga-explained:page-0071}\uses{}
\begin{verbatim}
A fiberedcategoryT overC isrepresentableifand only if
T isfiberedin groupoids,and thereisan object U of C and an object^ of ^{U),
such that forany objectp of ^ thereexistsa unique arrow p ?> ^ in ^.
    3.6.3. Splitting a fibered category. As we have seen in Example 3.14,a
fiberedcategorydoes not necessarilyadmit a splitting.
                                                    However, a fiberedcategory
isalways equivalentto a splitfiberedcategory.
\end{verbatim}
\end{proposition}
\begin{proof}
\notready
\end{proof}


\begin{proposition}\label{fga-pro-3-49}\dcref{3.49}\source{fga-explained:page-0076}\uses{}
\begin{verbatim}
Let p be an object of J^{X). To givep the structureof a
G-equivariantobject isthe same as assigningan isomorphism 0: prjp = a*p in
J^{G X X), such that the diagram
                                 (mcXidx
                         WlP



                                    E*p
commutes.
    When applied to the fiberedcategory of sheaves of some kind (forexample,
quasi-coherentsheaves) one gets preciselythe usual definitionof an equivariant
sheaf.
CHAPTER     4


                                  Stacks


               4.1. Descent of objects of fibered categories
    4.1.1. Gluing continuous maps and topological spaces. The following
isthe archetypalexample of descent.Take (Cont) to be the categoryof continuous
maps (thatis,the categoryofarrows in (Top),as in Example 3.15);thiscategoryis
fiberedover (Top) viathe functorP(cont)? (Cont) ?> (Top) sending each continuous
map to itscodomain. Now, suppose that /,:X -^ U and g: Y -^ U are two objects
of (Cont)mapping to the same objectU in (Top);we want to constructa continuous
map 0: X ?> y over U, that is,an arrow in (Cont)([/)= {Top/U). Suppose that
we are given an open covering{Ui} of [/,and continuousmaps 0^: f~^Ui -^ g~^Ui
over Ui] assume furthermorethat the restriction of (j)i      to f~^{Ui C\Uj) ?>
                                                      and (j)j
g~^{Ui n Uj) coincide.Then thereisa unique continuousmap cj):    X -^ Y over U
whose restrictionto each f~^Ui coincideswith /^.
    This can be writtenas follows.The category(Cont) isfiberedover (Top),and
if/: V ?> [/ is a continuous map, X ?> [/ an object of (Cont)([/)= (Top/[/),
then a pullback of X ?> [/ to V is given by the projectionV Xjj X ?> V. The
functor/*: (Cont)([/)-> (Cont)(y) sends each object X -^ U to V Xu X -^ V,
and each arrow in (Top/[/),givenby continuousfunction(j):  X -^Y over [/,to the
continuousfunction/*0 = idy Xjjf: V Xjj X -^ V XjjY.
    Suppose that we are given two topologicalspaces X and Y with continuous
maps X ?> 5 and Y -^ S. Consider the functor
                       Homc(X,y):    (Top/5) ?> (Set)
from the categoryof topologicalspaces over 5, definedin Section3.7.This sends
each arrow [/ ?> 5 to the setof continuousmaps Homt/(t/ Xs X,U x^ y) over U.
The actionon arrows isobtainedas follows:given a continuousfunctionf:V-^U,
we send each continuousfunction(/>: UxsX-^UxsY        to the function
   /*0 = idy X(l):VxsX     = V xu {U XsX)    ?>V   xu {U xsY) = V xs Y.
    Then the factthat continuousfunctionscan be constructedlocallyand then
glued togethercan be expressedby saying that the functor
                      Homc(X,y): (Top/5)?P ?> (Set)
isa sheafin the classical
                        topology of (Top).
    But thereismore: not only we can constructcontinuousfunctionslocally:we
can alsodo thisforspaces,although thisismore complicated.
\end{verbatim}
\end{proposition}
\begin{proof}
\notready
\end{proof}
